\chapter{PROJECT GENERAL CONTEXT AND PRESENTATION OF THE COMPANY}

\section{Introduction}

This chapter is organized to provide a clear overview of the academic and industrial
framework in which our graduation project was carried out. First, we present the academic
institution and the host company that supported the development of this work. Next, we
describe the organizational structure and services of the company, along with its main
products, in order to highlight the environment in which the project was conceived.
Finally, we outline the specifications of the project, including its objectives and the
concept of the General Purpose Input Output Box, which represents the core of our study.


\section{Host Organization}
This section presents the hosting organization, Be Wireless Solutions, and its role in
supporting innovative technological projects.

\subsection{Presentation of the Host Organization}

Be Wireless Solutions (BWS) is a Tunisian company founded in 2016, recognized as one of
the pioneers of Internet of Things (IoT) technologies in the region. In 2017, the company
obtained its official IoT license, and by 2019 it was awarded the national start‑up label,
confirming its innovative positioning in the Tunisian technology ecosystem. From its
inception, BWS has focused on designing and manufacturing IoT‑related technologies,
including smart sensors, intelligent boxes, and integrated platforms that allow clients to
monitor and control their systems in real time through web and mobile applications. What
distinguishes BWS is its ability to deliver complete solutions that combine hardware,
software, and communication layers, while also producing its own devices locally in
Tunisia.

The company has achieved remarkable growth in the domestic market, marked by a
deployment rate of connected devices that increased by more than 200\% compared to
competitors. This success has been recognized internationally: in 2020, the American
magazine *Startup City* listed BWS among the ten most promising technology start‑ups in
Africa. In 2021, the company won the award for Best Small and Medium Enterprise (SME)
during the Tunisia Digital Days at Expo 2020 in Dubai. In 2022, BWS reached the finals of
the MED’INNOVANT AFRICA competition and was named by *MEA Markets* as the best IoT
company in Tunisia.

Beyond awards, BWS has built its reputation on sustainability, innovation, and customer
satisfaction. The company operates entirely from Tunisia, with its own production facilities
dedicated to IoT devices and solutions. Its expertise covers resource optimization,
including reducing fuel and maintenance costs, improving energy efficiency in industrial
and agricultural sectors, and enabling smart management of fleets, public lighting, and
critical infrastructure. BWS solutions also address environmental protection, monitoring
electricity consumption, water usage, cold chain logistics, and remote machine control,
providing clients with integrated dashboards, alert systems, and AI‑based recommendations
to support decision‑making.

Thanks to its capacity to deliver reliable and innovative solutions, BWS has expanded its
reach beyond Tunisia, establishing collaborations in France, Switzerland, Morocco, Senegal,
and Qatar. These international partnerships reinforce its role as a trusted provider of IoT
solutions at a global scale. Ultimately, the company’s strategic vision is to offer creative
and sustainable technologies that preserve resources while enabling organizations to
achieve greater efficiency and resilience in a connected world.

(image de logo BWS)

\subsection{Organization Chart and Company Services}
This subsection details the company’s departments and services, explaining how each
contributes to the overall mission of delivering technological solutions.
    \begin{figure}[H]
        \centering
        \includegraphics[width=0.7\textwidth]{C:/Users/selim/Downloads/imagesreport/organigram.png}
        \caption{Organization Chart}
        \label{fig:electricity}
    \end{figure}

\subsection{Projects and Services of Be Wireless Solutions}

Be Wireless Solutions provides several innovative IoT services and solutions. The main
projects and services include:

\begin{itemize}
    \item \textbf{Electricity Consumption Monitoring}: A smart energy management solution
    that identifies energy‑consuming equipment, provides real‑time monitoring, and offers
    AI‑based recommendations to reduce costs and improve carbon footprint.
    \begin{figure}[H]
        \centering
        
\includegraphics[width=0.7\textwidth]{C:/Users/selim/Downloads/imagesreport/electricity.png}
        \caption{Electricity consumption monitoring solution}
        \label{fig:electricity}
    \end{figure}

    \item \textbf{Water Consumption Monitoring}: A system that detects leaks and fraud,
    enables remote reading of water meters, and helps optimize water usage without
    requiring major infrastructure changes.
    \begin{figure}[H]
        \centering
        
\includegraphics[width=0.7\textwidth]{C:/Users/selim/Downloads/imagesreport/water.png}
        \caption{Water consumption monitoring solution}
        \label{fig:water}
    \end{figure}

    \item \textbf{Fleet and Fuel Management}: An integrated vehicle tracking solution that
    monitors fuel consumption in real time, supports eco‑driving practices, and improves
    operational efficiency for transport and logistics.
    \begin{figure}[H]
        \centering
        
\includegraphics[width=0.7\textwidth]{C:/Users/selim/Downloads/imagesreport/fleet.png}
        \caption{Fleet and fuel management solution}
        \label{fig:fleet}
    \end{figure}

    \item \textbf{Remote Control of Equipment}: A cloud‑enabled system that allows
    operators to remotely control machines, detect anomalies, and predict breakdowns,
    reducing waste and improving reliability in industrial processes.
    \begin{figure}[H]
        \centering
        
\includegraphics[width=0.7\textwidth]{C:/Users/selim/Downloads/imagesreport/equipment.png}
        \caption{Remote control of equipment solution}
        \label{fig:equipment}
    \end{figure}

    \item \textbf{Cold Chain Monitoring}: A wireless monitoring solution that ensures
    compliance with cold chain requirements in warehouses and transport vehicles, with
    instant alerts when temperature thresholds are exceeded.
    \begin{figure}[H]
        \centering
        
\includegraphics[width=0.7\textwidth]{C:/Users/selim/Downloads/imagesreport/coldchain.png}
        \caption{Cold chain monitoring solution}
        \label{fig:coldchain}
    \end{figure}
\end{itemize}


\section{Project Specification}
This section defines the specifications of the IO-Box project, outlining its scope and
requirements.

\subsection{Problematics}

One of the main challenges encountered by Be Wireless Solutions lies in the complexity and
time required to configure and deploy sensors in industrial environments. Each new sensor
installation traditionally involves a lengthy process of programming, wiring, and adapting
interfaces to ensure compatibility with existing systems. This repetitive and resource‑intensive
workflow not only slows down deployment but also increases the risk of errors and
maintenance difficulties.

To overcome these limitations, BWS has proposed the development of the IO‑Box, a
general‑purpose interfacing platform designed to simplify sensor integration.

\subsection{General Purpose Input Output Box}

A General Purpose Input Output (GPIO) Box is a modular device designed to provide a
flexible interface between digital systems and external hardware. It acts as a bridge that
allows microcontrollers, processors, or embedded platforms to interact with sensors,
actuators, and other peripherals in a standardized way. The GPIO Box typically offers a
set of configurable pins that can be programmed either as inputs, to read signals from
external devices, or as outputs, to control or drive external components.

In practice, a GPIO Box is widely used in industrial automation, embedded systems, and
educational environments because of its versatility. It enables the acquisition of digital
and analog signals, supports basic control tasks, and facilitates communication with
external modules. By abstracting the complexity of direct hardware connections, the GPIO
Box provides a simple and reliable interface for experimentation, prototyping, and system
integration.

The main advantages of a General Purpose Input Output Box are its adaptability and ease
of use. It can be reconfigured to meet different application requirements, ranging from
basic signal monitoring to more advanced control operations. This makes it a fundamental
building block in modern embedded system design, offering engineers and developers a
practical tool for connecting and managing diverse hardware components.

\subsection{Methodology: V-Model}

For the development of the IO‑Box project, the V‑Model (Verification and Validation
Model) was adopted as the methodological framework. This structured approach is
widely used in software engineering and embedded systems, particularly when a high
level of reliability is required. It is based on a rigorous sequential process that organizes
development and testing activities into two symmetrical branches forming a "V"
(Figure~\ref{fig:vmodel}).

On the descending branch, the process begins with the specification of requirements,
followed by functional design and detailed design. These steps ensure that the system’s
requirements are well defined and that a coherent hardware and software architecture
is established. Once the design is completed, the implementation (coding) phase
constitutes the core of the technical realization.

The ascending branch is dedicated to verification and validation, including unit tests,
integration tests, system tests, and final validation. Each of these phases corresponds
to a previous design stage, ensuring that every requirement is verified by appropriate
tests. This symmetry between design and testing promotes complete traceability of
requirements, which is a key factor in critical projects.

One of the main advantages of the V‑Model is the early detection of anomalies during
the initial phases of development, significantly reducing the risk of costly errors at the
end of the cycle. Moreover, this approach strengthens the quality of the final product
by imposing rigorous validation at each level and encouraging clear documentation
throughout the project. The V‑Model was therefore particularly suited to the context of
this project, which required a high degree of reliability and strict management of the
development cycle.

The V‑Model process applied to this project included the following steps:
\begin{itemize}
    \item \textbf{System Requirements}: Define and document the functional and
    technical requirements of the IO‑Box, including accuracy of measurements,
    real‑time constraints, and reliability of communication.
    \item \textbf{System Design}: Establish the overall architecture of the system,
    including sensors (ENS210, OPT3001, current sensor), STM32 microcontroller,
    and communication interfaces.
    \item \textbf{Software Requirements}: Specify the software modules required,
    such as signal acquisition, filtering, regulation, and communication protocols.
    \item \textbf{Software Implementation}: Develop the firmware for STM32 using
    IAR Embedded Workbench, ensuring proper configuration of drivers, BSP, and HAL.
    \item \textbf{Software Integration}: Integrate communication, signal processing,
    and control modules into a coherent application.
    \item \textbf{System Integration}: Deploy the firmware on the STM32 board and
    connect the system to sensors and actuators for autonomous operation.
    \item \textbf{Validation Tests}: Perform functional and performance tests to verify
    measurement accuracy, stability of regulation, and reliability of communication
    with external supervisory systems.
\end{itemize}

\begin{figure}[H]
    \centering
    \includegraphics[width=0.75\textwidth]{C:/Users/selim/Downloads/imagesreport/v-model.png}
    \caption{Development cycle based on the V‑Model}
    \label{fig:vmodel}
\end{figure}


\subsection{Objectives}

The main objective of the IO‑Box project is to design and implement a versatile
interfacing and signal‑processing platform capable of acquiring, processing, and
communicating industrial data in real time. More specifically, the system aims to:

\begin{itemize}
    \item \textbf{Signal Acquisition}: Collect heterogeneous analog and digital signals
    from various sensors and industrial equipment.
    
    \item \textbf{Interfacing with External Devices}: Communicate with third‑party sensors
    and actuators using multiple industrial protocols.
    
    \item \textbf{Local Data Processing}: Perform on‑board filtering, calculations, and
    regulation to ensure reliable and efficient operation.
    
    \item \textbf{Actuator Control}: Drive actuators with precision and robustness,
    guaranteeing dependable system responses.
    
    \item \textbf{External Communication}: Exchange data with supervisory systems,
    applications, or IoT platforms for monitoring and integration.
\end{itemize}

On the software side, the project emphasizes optimal exploitation of hardware resources
through a modular architecture based on hardware abstraction layers (HAL). This includes
driver management for peripherals and independence from the underlying hardware.

\subsection{Expected Impact}

The IO‑Box is expected to simplify sensor integration, reduce deployment time, and
improve operational efficiency in industrial contexts. By offering a standardized
interface, it enhances interoperability between devices and ensures scalability for future
applications. This impact positions the IO‑Box as a practical contribution to modern
industrial IoT solutions.

\subsection{Conclusion}

In this chapter, we presented the general framework of the project. We began with an
overview of Be Wireless Solutions and highlighted the main problematics related to sensor
integration. This led us to define the objectives of the IO‑Box project, which aims to
provide a modular and efficient platform for signal acquisition, processing, and
communication. By clarifying the context and goals, this chapter sets the foundation for
the detailed study that follows.

The next chapter will be devoted to the analysis and collection of project requirements.
It will specify the functional and technical needs that the IO‑Box must address, serving as
a basis for the design and implementation phases.
